\documentclass[11pt,letterpaper]{article}
\usepackage[margin=0.8in]{geometry}
\usepackage[T1]{fontenc}
\usepackage{lmodern}
\usepackage{amsmath,amssymb,xcolor,listings,enumitem,hyperref}
\hypersetup{colorlinks=true,linkcolor=blue,urlcolor=blue,plainpages=false,pdftitle={CMSC433 Topic Practice}}
\definecolor{paper}{RGB}{246,248,250}
\lstset{basicstyle=\ttfamily\small,backgroundcolor=\color{paper},columns=fullflexible,
keepspaces=true,breaklines=true,showstringspaces=false,frame=single,
rulecolor=\color{gray!30},aboveskip=0.5em,belowskip=0.5em}
\lstnewenvironment{codeanswer}{\lstset{backgroundcolor=\color{white},frame=single}}{}
\setlist[enumerate]{itemsep=0.7em}
\setlength{\parindent}{0pt}
\setlength{\parskip}{0.5em}
\ExplSyntaxOn
\NewDocumentCommand{\code}{m}{\texttt{%
\tl_set:Nx\l_tmpa_tl{\tl_to_str:n{#1}}%
\seq_set_split:NnV\l_tmpa_seq{~}\l_tmpa_tl
\seq_map_indexed_inline:Nn\l_tmpa_seq{%
\int_compare:nNnT{##1}>{1}{\penalty0\relax\space}%
\str_map_inline:nn{##2}{####1\kern0pt\relax}}}}
\ExplSyntaxOff
% TYPE YOUR WRITTEN ANSWERS INSIDE THESE BRACES in the included lesson.tex files.
\newcommand{\answer}[1]{\par\nopagebreak[4]\smallskip\noindent\textbf{Answer: }#1\par\medskip}
\newcommand{\task}[1]{\subsection*{#1}}
\newcommand{\source}[1]{{\small\textbf{Source:} #1}\par}
\newcommand{\ready}[1]{{\small\textbf{Checkpoint:} #1}\par}

\raggedbottom
\newcommand{\sheettitle}[3]{\noindent{\small CMSC433 / Part #1}\par
{\Large\bfseries #2}\par{\large #3}\par\smallskip\hrule\medskip}
\newcommand{\mockrules}{{\small Topic mock, using the sample's question styles.
Attempt independently; feedback after the whole file. Untimed practice, no points
assigned. Current exam scope and rules are unconfirmed.}\par}

\begin{document}
\sheettitle{06}{Verification conditions}{Guided practice}

\textbf{Use with the chat:} do Question 1 first; then derive one loop obligation
at a time. These are partial-correctness questions; omit termination obligations.

\textbf{Worked example.} For \code{x := x - 2;} with postcondition $x\geq4$,
substitution gives $x-2\geq4$, equivalently $x\geq6$. Work backward in a sequence.

\task{1. Weakest precondition}
Compute the weakest precondition of \code{x := y + 2; y := x - 1;}
for postcondition $x>y\land y\geq0$. Show the intermediate formula.
\answer{Before last assignment: [ ]; before first: [ ].}

\task{2. One invariant, three obligations}
\par\noindent\begin{minipage}{\linewidth}
\begin{lstlisting}
method Twice(n:int) returns (t:int)
  requires n >= 0
  ensures t == 2*n + 1
{
  t := 0;
  var i := n;
  while i > 0 {
    t := t + 2;
    i := i - 1;
  }
  t := t + 1;
}
\end{lstlisting}
\end{minipage}\par
Use invariant $I\equiv i\geq0\land t+2i=2n$.
Derive initialization, preservation, and exit, substituting the code into each.
Account for the final assignment when deriving the exit obligation.
\answer{Initialization: [ ].}
\answer{Preservation: [ ].}
\answer{Exit: [ ].}
\source{Sample Q3, p. 3; DafnyVerificationConditions physical pp. 3--30.}

\end{document}
